% Chapter 6 -- Rubric: PO11, PO9 (9.1--9.4), PO6 (6.1), PO7 (7.1), PO8 (8.1--8.4)
\chapter{Project Management, Teamwork, and Impact}
\label{ch:management}

\section{Project Management and Finance}

\subsection{Planning and Risk Management}
\label{sec:plan}

The project followed an agile, iterative research and development lifecycle organized into six primary chronological phases:
\begin{enumerate}
  \item \textbf{Phase 1: Problem Conception and Literature Review} \pending{Reconstructed dates: January--February 2026}. Established conversational asymmetry foundations, reviewed P4G literature, and formulated core problem definitions.
  \item \textbf{Phase 2: Intention Ground-Truth Curation} \pending{Reconstructed dates: March--April 2026}. Conducted 5-way pilot calibration, manual annotation of 1,017 P4G dialogues (Pass 1), and inter-annotator drift verification.
  \item \textbf{Phase 3: Intention Model Exploration (Stages 1--3)} \pending{Reconstructed dates: May--July 2026}. Developed flat models (v1--v5), diagnosed v3 metric bugs, implemented hierarchical dialogue transformers (v6--v9), and investigated last-turn gating and routing (v10--v12).
  \item \textbf{Phase 4: Strategy Taxonomy and Multi-Label Annotation} \pending{Reconstructed dates: July--August 2026}. Defined the 11-category, 41-strategy hierarchy, developed the contextual batching pipeline, and annotated $10{,}600$ Persuader turns under human checkpoints.
  \item \textbf{Phase 5: Model Consolidation and E-Commerce Benchmarking} (Confirmed Git history: September 5--15, 2026). Deprecated modifier head (Pass 2), developed Antiphony (v14), synthesized the 700-dialogue negotiation benchmark, executed strategy ensemble searches (v17/v18), and conducted adversarial ablations.
  \item \textbf{Phase 6: Strict Re-annotation and Thesis Finalization} (Confirmed Git history: September 16--20, 2026). Executed strict third-pass intention relabeling (694/323), completed skew stress runs, performed forensic post-mortems, and compiled this report.
\end{enumerate}

\begin{figure}[htbp]
  \centering
  \begin{ganttchart}[
      hgrid, vgrid, x unit=0.95cm,
      time slot format=isodate-yearmonth, time slot unit=month,
      title label font=\small\bfseries, bar label font=\footnotesize,
      bar/.style={fill=DecMain!40, draw=DecBord},
      bar incomplete/.style={fill=PerMain!40, draw=PerBord}
    ]{2026-01}{2026-10}
    \gantttitlecalendar{year, month=shortname} \\
    \ganttbar{1. Conception \& Literature \pending{Est.}}{2026-01}{2026-02} \\
    \ganttbar{2. Intention Curation (Pass 1) \pending{Est.}}{2026-03}{2026-04} \\
    \ganttbar{3. Architectural Lineage (v1--v12) \pending{Est.}}{2026-05}{2026-07} \\
    \ganttbar{4. Strategy Annotation (10.6k) \pending{Est.}}{2026-07}{2026-08} \\
    \ganttbar{5. Antiphony \& Datagen (v13/v14)}{2026-09}{2026-09} \\
    \ganttbar{6. Pass 3 \& Thesis Finalization}{2026-09}{2026-10}
  \end{ganttchart}
  \caption{Project timeline and phase distribution. Calendar dates prior to September 5, 2026 are reconstructed estimates (\pending{Confirmed Git history begins September 5, 2026}).}
  \label{fig:gantt}
\end{figure}

\Cref{tab:risk-matrix} presents the proactive risk management matrix implemented during system development.

\begin{table}[htbp]
  \centering
  \small
  \caption{Project risk management matrix, impact levels, and mitigations}
  \label{tab:risk-matrix}
  \begin{tabularx}{\textwidth}{>{\raggedright\arraybackslash\bfseries}p{0.24\textwidth} >{\raggedright\arraybackslash}p{0.18\textwidth} cc >{\raggedright\arraybackslash}X}
    \toprule
    \textbf{Risk Event} & \textbf{Risk Category} & \textbf{Prob.} & \textbf{Impact} & \textbf{Proactive Mitigation Strategy} \\
    \midrule
    Ground-Truth Label Noise in P4G & Data Quality & High & High & Executed 3-pass manual re-annotation, establishing strict unconditional commitment ground truth. \\
    Extreme Long-Tail Strategy Sparsity & Modeling Feasibility & High & High & Pivoted primary target to 11 communicative categories; formulated multi-task Asymmetric Loss regularizer. \\
    Hardware GPU Memory (OOM) on Hierarchical Models & Technical Constraints & Med & High & Implemented micro-batch slicing ($B=2$), gradient checkpointing, and dynamic memory allocation guards. \\
    Subjective Drift in Strategy Annotation & Data Integrity & Med & Med & Enforced 4-tier decision protocol (Removal Test, Density Cap) and automated closed-world validation checks. \\
    Tokenizer Truncation Artifacts & Evaluation Validity & Med & High & Conducted forensic token length audits; converted default right-truncation to left-truncation. \\
    \bottomrule
  \end{tabularx}
\end{table}

\subsection{Budgeting and Resource Identification}
\label{sec:budget}

To demonstrate that publication-grade conversational AI research can be achieved responsibly, the project operated under strict zero-budget computational constraints:
\begin{itemize}
  \item \textbf{Human Capital:} Five undergraduate student researchers providing domain research, manual dialogue annotation, prompt engineering, model implementation, and empirical evaluation.
  \item \textbf{Compute Hardware:} Exclusively utilized free-tier cloud accelerators: Kaggle Jupyter environments equipped with a single NVIDIA Tesla T4 GPU (16\,GB GDDR6 VRAM), 2 vCPUs, and 13\,GB system RAM. Total zero expenditure on dedicated institutional compute clusters.
  \item \textbf{Software and Pre-trained Checkpoints:} Open-source libraries (PyTorch, Hugging Face, Scikit-learn, XGBoost) and publicly hosted pre-trained checkpoints (RoBERTa, DeBERTa-v3, TOD-BERT, ELECTRA) under permissive research licenses.
  \item \textbf{API Services for Synthetic Generation:} Paraphrasing and semantic verification leveraged free-tier and low-cost developer promotional credits for the Google Gemini API \pending{Total commercial API expenditure: under \$15, confirm exact billing ledger}.
\end{itemize}

\subsection{Economic Analysis and Financial Prospecting}
\label{sec:economics}

In emerging conversational commerce ecosystems, micro-merchants experience severe operational inefficiency, spending up to 60\% of their daily labor answering non-converting inquiries. The economic value proposition of deploying compact, discriminative intention classifiers includes:
\begin{enumerate}
  \item \textbf{Lead Triage Automation:} By continuously scoring purchasing probability ($\hat{y}$), customer support systems can automatically route high-intent buyers to senior sales representatives while serving automated informational responses to casual exploratory browsers.
  \item \textbf{Low Operating Costs:} Because Antiphony runs on standard commodity hardware ($<50$\,ms latency on single GPUs without requiring billion-parameter LLM clusters), operational deployment costs are orders of magnitude lower than querying proprietary commercial LLM APIs (\$0.00002 per turn vs. \$0.02+ per API call).
  \item \textbf{Sales Coaching and Quality Assurance:} Quantifying persuader strategy distributions across sales teams identifies effective rhetorical patterns, improving staff training in non-profit fundraising and retail commerce.
\end{enumerate}

\subsection{Sustainability in Business and Commercialisation}
\label{sec:commercialisation}

To remain commercially viable post-course, the framework can be adapted as an open-core software library or integrated into social commerce CRM platforms (e.g., Messenger management dashboards). Maintaining operational viability requires: (i) periodic retraining as colloquial slang and product catalogs shift, (ii) establishing data feedback loops where sales reps flag false positives, and (iii) containerizing the inference pipeline into lightweight Docker microservices.

\section{Individual and Teamwork}
\label{sec:team}

\subsection{Participation and Contribution}
\label{sec:participation}

The capstone project was executed by five undergraduate students in the Department of Computer Science and Engineering, BUET (Group 10): **Arnob Biswas, Nawriz Ahmed Turjo, Abhishek Roy, Monjur Hossain Khan, and Shams Hossain Simanto**. 

Contributions documented directly in the project repository and research chronicle include:
\begin{itemize}
  \item \textbf{Dataset Annotation and Ground Truth Curation:} All five team members actively participated in the five-way pilot calibration and independently hand-annotated approximately 200 dialogues each to construct the primary 1,017-dialogue P4G ground-truth corpus. Team members also conducted the independent second-pass agreement study (154 dialogues) and the terminal stance verification of the 81 rechecked dialogues.
  \item \textbf{Strategy Annotation and Validation:} The four strategy annotation segments were conducted under human checkpoints: \texttt{shimanto} (turns 1--2000), \texttt{shovon} (turns 2001--4000, 6001--7000), unassigned batch (turns 4001--6000), and \texttt{abhishek} (turns 7001--10600).
  \item \textbf{Synthetic E-Commerce Benchmark Audit:} All team members participated in the exhaustive 100\% manual verification of the 700 synthetic negotiation dialogues, verifying conversational consistency, correcting label drift, and validating reversal turn markers.
  \item \textbf{Codebase and Engineering Contributions (Git Log: Sept 5--20, 2026):}
    \begin{itemize}
      \item \textbf{Abhishek Roy (9 commits):} Orchestrated the Antiphony dual-stream architecture, multi-task strategy modeling (v17/v18), decision combiner searches, and cross-domain adversarial ablation pipelines.
      \item \textbf{Nawriz Ahmed Turjo (3 commits):} Managed dataset structuring, Pass 3 deterministic remapping, token-safe augmentation scripts, and baseline benchmarking.
      \item \textbf{Shams Hossain Simanto (2 commits):} Coordinated synthetic datagen humanization pipelines, exploratory data analysis scripts, and visualization generation.
      \item \textbf{Arnob Biswas and Monjur Hossain Khan:} \pending{Confirm specific individual engineering module allocations beyond manual annotation and joint auditing}.
    \end{itemize}
\end{itemize}

\subsection{Collaboration and Conflict Resolution}
\label{sec:collaboration}

Collaboration was structured through shared cloud repositories, standardized Jupyter notebook conventions, and bi-weekly synchronization meetings. When empirical conflicts arose—such as the collapse of v3 training or disagreements on boundary intention labels—the team resolved them through objective, data-driven protocols:
\begin{itemize}
  \item \textit{Metric Adjudication:} When v3 collapsed, rather than discarding the run or guessing causes, the team systematically audited checkpoint logs, uncovering the metric bug and class weight explosion.
  \item \textit{Consensus Labeling:} Disagreements on ambiguous commitment dialogues were resolved through joint panel re-reads, establishing strict, unambiguous rules (e.g., classifying all temporal deferrals as \texttt{no}).
\end{itemize}

\subsection{Leadership and Direction}
\label{sec:leadership}

Technical direction was steered through clear task division and proactive recovery planning. Following the supervisor's critical directives—specifically the mandate to drop the modifier head and tighten the commitment definition—the team demonstrated agility by halting unproductive multi-task branches and systematically refactoring the entire evaluation harness within 72 hours.

\subsection{Multidisciplinary Engagement}
\label{sec:multidisciplinary}

The project engaged deeply across multiple academic disciplines:
\begin{itemize}
  \item \textbf{Behavioral Psychology and Rhetoric:} Translating Cialdini's influence theory \cite{cialdini2001influence} and sequential compliance tactics \cite{freedman1966compliance,cialdini1975reciprocal} into a formalized computational taxonomy.
  \item \textbf{Computational Linguistics and Pragmatics:} Analyzing conversational implicature, speech acts, and turn-taking discourse dynamics.
  \item \textbf{E-Commerce and Emerging Market Economics:} Grounding adversarial benchmark scenarios in the practical realities of mobile social commerce in Bangladesh.
\end{itemize}

\section{Impact on Society}
\label{sec:society}

\subsection{Positive Societal and Commercial Benefits}
The developed computational models offer substantial positive social utility:
\begin{itemize}
  \item \textbf{Empowerment of Independent MSMEs:} Providing resource-constrained micro-merchants with automated customer triage tools, allowing them to compete effectively against large commercial retail platforms.
  \item \textbf{Enhancing Non-Profit Philanthropy:} Providing charitable organizations with empirical insights into which communicative strategies foster genuine empathy and donor engagement, improving fundraising efficiency for humanitarian crises.
  \item \textbf{Customer Autonomy and Transparency:} Automated strategy classification tools can be inverted to act as consumer protection guards, alerting online shoppers when conversational agents deploy aggressive scarcity pressure, artificial deadlines, or guilt-inducing rhetoric.
\end{itemize}

\subsection{Dual-Use Risks: Coercive Manipulation and Dark Patterns}
Like all persuasive technologies, computational persuasion models carry inherent dual-use risks:
\begin{itemize}
  \item \textbf{Automated Exploitative Persuasion:} Malicious actors could leverage strategy optimization models to design predatory chatbots that systematically exploit psychological vulnerabilities (e.g., targeting low-income or elderly users with coercive debt-financing or deceptive gambling products).
  \item \textbf{Mitigation:} In accordance with IEEE 7000 value-based design, this project explicitly restricted its scope to \textit{discriminative analysis and understanding}, deliberately excluding natural language generation (NLG) engines capable of autonomous persuasive manipulation.
\end{itemize}

\section{Environmental Impact and Life Cycle Sustainability}
\label{sec:sustainability}

\subsection{Transparent Computational Energy and Carbon Accounting}
In accordance with modern sustainable computing mandates \cite{strubell2019energy,lacoste2019quantifying}, we transparently account for the environmental footprint of our experimental research. All training runs were conducted on cloud-hosted NVIDIA Tesla T4 GPUs (rated at a maximum thermal design power of 70\,W board draw). 

\Cref{tab:carbon-footprint} itemizes the estimated compute hours across the reported experimental lineages.

\begin{table}[htbp]
  \centering
  \small
  \caption{Estimated computational energy consumption and greenhouse gas emissions}
  \label{tab:carbon-footprint}
  \begin{tabularx}{\textwidth}{l>{\raggedleft\arraybackslash}X>{\raggedleft\arraybackslash}X>{\raggedleft\arraybackslash}X}
    \toprule
    \textbf{Experimental Lineage / Phase} & \textbf{GPU Hours} & \textbf{Energy (kWh)} & \textbf{Emissions (kg CO$_2$e)} \\
    \midrule
    Intention Exploration (Stages 1--3: v1--v12) & 24.5 h & 1.715 kWh & 0.86 -- 1.20 kg \\
    Antiphony Development \& Pass 2/3 (v13--v14) & 16.0 h & 1.120 kWh & 0.56 -- 0.78 kg \\
    Strategy Classifier Reruns (Clean + v17 + v18) & 12.0 h & 0.840 kWh & 0.42 -- 0.59 kg \\
    E-Commerce Adversarial Transfer \& Ablations & 8.5 h & 0.595 kWh & 0.30 -- 0.42 kg \\
    \midrule
    \textbf{Total Estimated Research Compute} & \textbf{61.0 h} & \textbf{4.270 kWh} & \textbf{2.14 -- 2.99 kg CO$_2$e} \\
    \bottomrule
  \end{tabularx}
  \vspace{2pt}
  \footnotesize{\newline *Note: Energy computed assuming continuous 70\,W GPU draw; carbon emissions estimated using regional grid carbon intensity factors of $0.50$--$0.70$\,kg CO$_2$e/kWh \cite{lacoste2019quantifying}. Excludes unlogged exploratory runs and CPU data preparation.}
\end{table}

The total estimated carbon footprint of $\approx 2.5$\,kg CO$_2$e is modest (comparable to driving an average passenger vehicle for less than 10 kilometers). We note that earlier unverified prompts cited an unsubstantiated estimate exceeding eighteen kilograms; our transparent, hardware-grounded derivation demonstrates the true order of magnitude.

\section{Ethics}
\label{sec:ethics}

\subsection{Equity and Inclusivity}
\label{sec:equity}
The project acknowledges significant linguistic and cultural boundaries:
\begin{itemize}
  \item \textbf{Linguistic Exclusion:} All investigated datasets and models operate exclusively in English. This excludes non-English speakers and fails to capture the code-mixed dialects (e.g., Banglish) prevalent in South Asian social commerce.
  \item \textbf{Cultural Specificity:} The P4G corpus reflects Western, North American conversational norms regarding charitable giving. Applying these rhetorical taxonomies to non-Western contexts without cultural adaptation may introduce systematic evaluation bias.
\end{itemize}

\subsection{Accountability and Personal Responsibility}
\label{sec:accountability}
The authors accept full personal and academic responsibility for the integrity, methodology, and reported metrics of this report. All reported failures (e.g., v3 collapse, tokenizer truncation artifacts, and the rejection of the XGBoost stacker) are disclosed transparently to uphold scientific accountability.

\subsection{Proper Use of Intellectual Property}
\label{sec:ip}
All third-party assets, pre-trained model checkpoints, and open-source packages are used in full compliance with their respective academic and open-source licenses:
\begin{itemize}
  \item \textbf{Datasets:} Persuasion for Good (P4G) is utilized under its standard academic research license \pending{Confirm license details}.
  \item \textbf{Pre-Trained Encoders:} RoBERTa, DeBERTa-v3, TOD-BERT, and ELECTRA are distributed under Apache 2.0 and MIT open-source licenses.
  \item \textbf{Disclosure of Generative AI Assistance:} In accordance with BUET academic integrity standards, we explicitly disclose the use of generative AI technologies: (i) Claude-3.7/Opus agent pipelines were utilized for multi-label persuasion strategy annotation under human checkpoints, (ii) Google Gemini Flash API was employed for organic paraphrasing in synthetic e-commerce dialogue generation, and (iii) an agentic AI assistant was utilized as a typesetting and drafting aid in compiling this LaTeX thesis from the team's underlying research notes.
\end{itemize}

\subsection{Professionalism and Ethical Codes}
\label{sec:codes}
The project adheres to the highest standards of professional engineering ethics, conforming strictly to the ACM Code of Ethics \cite{acmcode} and the IEEE Code of Ethics \cite{ieeecode}, ensuring that computational capabilities advance human understanding while actively mitigating risks of technological harm.
